\documentclass{article} % For LaTeX2e
\usepackage{iclr2026_conference,times}

% NOTE: the ICLR template ships a \input{math_commands.tex} line here. It is
% deliberately omitted: none of its macros are used, and it is one more file
% that must be present for the document to build at all.

\usepackage{amsmath}
\usepackage{amssymb}
\usepackage{hyperref}
\usepackage{url}
\usepackage{booktabs}
\usepackage{graphicx}
\usepackage{multirow}
\usepackage{xcolor}

% Marks a slot whose experiment has not finished yet.
\newcommand{\TODO}[1]{\textcolor{red}{[\textsc{todo}: #1]}}
\newcommand{\PEND}{\textcolor{gray}{--}}

\title{When Does a Token Graph Actually Help?\\
Calibration Failures in Graph-Based Sentence Pooling}

\author{Anonymous authors\\
Paper under double-blind review}

\begin{document}

\maketitle

\begin{abstract}
Graph-based pooling methods build a latent token-similarity graph on top of a
frozen language model and refine token representations with a graph neural
network before aggregation. We study \textsc{Glot}
\citep{mantri2026glot}, a recent method of this family, and show that the way
its graph is constructed does not transfer across backbones. \textsc{Glot}
thresholds token cosine similarity at an \emph{absolute} value $\tau$ selected
on BERT. Because the scale of token cosine similarity varies enormously across
language models, the same $\tau$ that yields a $0.103$-density graph on BERT
yields a density of $0.991$ on RoBERTa, $0.599$ on ModernBERT and $0.741$ on
SmolLM2. For RoBERTa the tenth-percentile token cosine is $0.693$, above
\emph{every} $\tau$ in the published search space $\{0.1, 0.3, 0.6\}$: no
configuration that was searched produces a sparse graph, so the reported
RoBERTa numbers are obtained from an essentially complete graph in which the
graph neural network receives no selective relational structure. We propose a
two-part correction --- relative-density thresholding and median-based input
scaling --- and show it recovers $+5.44$ MCC on ModernBERT/CoLA, while a BERT
control arm confirms the correction is inert where it should be. With graphs
calibrated, we then ask whether \emph{hyperbolic} geometry improves the pooler,
evaluating three extensions (a Poincar\'e token graph, a hyperbolic readout and
a hyperbolic GNN) and their combinations under an equal-budget search over the
paper's own optimizer and architecture grid, with 15 seeds and paired testing.
The result separates two components that prior framing conflated. The
hyperbolic \emph{readout} (Stage B) is significantly harmful: all four arms
containing it lose $0.75$--$1.24$ Spearman on STS-B with $0$--$1$ of $15$ seeds
positive, surviving Bonferroni correction. The hyperbolic \emph{GNN} (Stage C)
is positive against the baseline on both tasks ($+1.42$ MCC on CoLA, $13/15$
seeds; $+0.20$ Spearman on STS-B, $12/15$) --- but against a control that
deletes the token graph entirely it is a null ($+0.58$ MCC, $8/7$ seeds,
sign-test $p = 1.00$), because most of its apparent gain is shared with that
control. The hyperbolic token graph (Stage A) is likewise a null once the
Euclidean baseline is given the same architecture budget. Across every
backbone, task and search budget we measured, deleting the token graph is
indistinguishable from --- and on CoLA better than --- the full method, despite
well-conditioned graph density, so the constructed graph carries no measurable
information.
We further show that the published \textsc{Glot} numbers are single-seed maxima
over a hyperparameter grid: at the authors' own seed and grid we exceed their
reported CoLA by $3.03$ MCC and their RTE by $0.36$, and their reported
$\tau$-sensitivity curve spans a range ($7.87$ on CoLA) comparable to the
variation we measure across seeds at \emph{fixed} $\tau$ ($5.46$).
Along the way we document two measurement artifacts that we believe affect the
wider literature: $\delta$-hyperbolicity computed on raw LLM activations is
dominated by attention-sink norm outliers rather than linguistic structure, and
tuning-stage maxima in this setting are inflated by up to $1.83$ points by
selection bias alone.
\end{abstract}

%-------------------------------------------------------------------------
\section{Introduction}
\label{sec:intro}

Reducing a language model's token-level outputs to a single vector is a
prerequisite for nearly every sentence-level task. Standard poolers --- mean,
max, \texttt{[CLS]} --- treat tokens as an unordered set and discard the
relational structure that self-attention computed. A recent line of work
therefore re-frames pooling as \emph{relational learning followed by
aggregation}: build a latent graph over tokens, refine token states with a
graph neural network (GNN), then aggregate. \textsc{Glot}
\citep{mantri2026glot} is a representative and strong instance, reporting
competitive GLUE and MTEB performance with $20\times$ fewer trainable
parameters than comparable methods.

The premise of this family is that the constructed graph carries information.
Our starting point was the natural next question --- whether a \emph{hyperbolic}
graph would carry more, given that hyperbolic space embeds tree-like structure
with low distortion and that linguistic structure is often argued to be
hierarchical. In the course of testing that hypothesis we found that the
premise itself does not hold uniformly: for several widely used backbones the
graph that \textsc{Glot} constructs is degenerate, and the GNN operates on what
is effectively a complete graph.

The cause is a scale assumption. \textsc{Glot} connects tokens $i,j$ when
$\cos(x_i, x_j) > \tau$ with $\tau$ an absolute threshold searched over
$\{0.1, 0.3, 0.6\}$. This is only meaningful if the distribution of token cosine
similarity is comparable across models. It is not. Anisotropy of contextual
embeddings varies widely between language models, and a threshold calibrated on
one backbone produces graphs of wildly different density on another
\TODO{cite anisotropy literature}.

\paragraph{Contributions.}
\begin{enumerate}
\item \textbf{A calibration failure in graph-based pooling}
      (\S\ref{sec:calibration}). We measure the edge density that
      \textsc{Glot}'s own published $\tau$ grid induces on five backbones. On
      BERT the grid spans a useful range ($0.850$--$0.149$); on RoBERTa every
      value in the grid gives density $\geq 0.992$. We show the mechanism ---
      RoBERTa's tenth-percentile token cosine ($0.701$) lies above the largest
      $\tau$ searched --- and discuss which published results this affects.
\item \textbf{A correction that works, with a control}
      (\S\ref{sec:fix}). Replacing the absolute threshold with a
      relative-density criterion, and rescaling token norms by their
      \emph{median} rather than mean, recovers $21.71 \rightarrow 27.15$ MCC on
      ModernBERT/CoLA. Neither correction alone suffices: density-matching by
      itself \emph{degrades} performance to $9.16$. A BERT control arm, which
      should not benefit, correctly does not.
\item \textbf{A well-powered negative result on hyperbolic pooling}
      (\S\ref{sec:hyperbolic}). Over two tasks, five backbones, 15 paired
      seeds and matched tuning budgets across the paper's full optimizer and
      architecture grid, the hyperbolic variants separate rather than
      uniformly fail. The hyperbolic \emph{readout} is significantly harmful
      ($0/15$ seeds positive, surviving Bonferroni); the hyperbolic \emph{GNN}
      is never harmful and directionally positive; the hyperbolic \emph{token
      graph} is a null once the Euclidean control is given the same
      architecture budget. Referenced against a no-graph control rather than
      against \textsc{Glot}'s baseline, none of them is a win.
\item \textbf{Evidence that the graph itself may be dispensable}
      (\S\ref{sec:vs-nograph}, \S\ref{sec:decoder}). On both BERT and
      TinyLlama-1.1B, deleting the token graph is statistically
      indistinguishable from the full cosine graph --- on CoLA it is
      \emph{better} ($+0.83$ MCC) --- even though those graphs are
      well-conditioned. This is a question about graph-based pooling in
      general, not about our extensions to it.
\item \textbf{The published numbers are grid maxima} (\S\ref{sec:estimator}).
      \textsc{Glot} fixes seed $42$ and reports the best $\tau$ per task on the
      split it also reports. At that seed and on that grid we \emph{exceed} the
      published CoLA ($+3.03$) and RTE ($+0.36$), while our 15-seed mean sits
      $2.12$ below the same CoLA number. We show the reported
      $\tau$-sensitivity curve is not separable from variation at fixed $\tau$.
\item \textbf{Two measurement artifacts} (\S\ref{sec:artifacts}). We show that
      $\delta$-hyperbolicity on raw LLM activations measures norm outliers, not
      hierarchy, and we quantify the selection bias that makes tuning-stage
      maxima unusable as evidence in this regime.
\end{enumerate}

\paragraph{What this paper is not.} We do not claim \textsc{Glot} is wrong. Its
BERT results replicate closely (\S\ref{sec:repro}) and its stress-test
robustness claim is unaffected by anything we report. Our claim is narrower and,
we think, more useful: the graph construction has an unstated calibration
dependency, and once that is fixed the case for adding geometry to the pooler
is weaker than it first appears.

%-------------------------------------------------------------------------
\section{Background}
\label{sec:background}

\subsection{Graph-based pooling}
Let $H = (x_1,\dots,x_L)$, $x_i \in \mathbb{R}^d$, be the token states of a
frozen language model. \textsc{Glot} constructs an adjacency
\begin{equation}
A_{ij} = \mathbf{1}\!\left[\, \cos(x_i, x_j) > \tau \,\right],
\label{eq:glot-adj}
\end{equation}
applies $K$ graph attention layers with jumping-knowledge concatenation, and
aggregates the refined states with a learned readout. Only the pooler is
trained; the backbone is frozen and its hidden states are precomputed.

\TODO{expand: exact GAT configuration, readout, JK variant, parameter count}

\subsection{Hyperbolic extensions under test}
\label{sec:stages}
We consider three orthogonal modifications, which we refer to as Stages A, B
and C, and all their combinations.

\paragraph{Stage A --- Poincar\'e token graph.} Token states are mapped into the
Poincar\'e ball $\mathbb{D}^d_c$ via $\exp^c_0(\cdot)$ after optional centering
and normalisation, and edges are formed by thresholding \emph{geodesic}
distance rather than cosine similarity.

\paragraph{Stage B --- hyperbolic readout.} Aggregation is performed as a
Fr\'echet-style mean in the Poincar\'e ball rather than in $\mathbb{R}^d$.

\paragraph{Stage C --- hyperbolic GNN.} Message passing is performed with
M\"obius operations in $\mathbb{D}^d_c$.

\TODO{formal definitions of each stage; equations for expmap/logmap, M\"obius
addition, and the readout}

\paragraph{The effective-radius parameter.} Stage A has a subtlety that turns
out to matter. The behaviour of the Poincar\'e graph depends not on the
curvature $c$ alone but on the \emph{effective radius}
$r_{\text{eff}} = \sqrt{c}\,\lVert x \rVert$ after embedding. When
$r_{\text{eff}}$ is small the geodesic ordering degenerates to the Euclidean
ordering; when it is large, $\exp^c_0$ saturates and the ordering degenerates
to a cosine ordering. We find a usable window of
$r_{\text{eff}} \in [0.8, 2.5]$ and treat it as a search dimension rather than
fixing $c = 1$. \TODO{figure: task metric and edge density vs $r_{\text{eff}}$}

%-------------------------------------------------------------------------
\section{Experimental protocol}
\label{sec:protocol}

Because several of our claims are negative, the protocol is stated before any
result. Each element below was adopted in response to a specific artifact we
observed; we report those artifacts in \S\ref{sec:artifacts}.

\begin{enumerate}
\item \textbf{Two-stage evaluation.} A \emph{tuning} stage draws $T$
      hyperparameter configurations per arm at a fixed seed; a
      \emph{confirmation} stage re-runs the selected configuration on 15
      held-out seeds. Only confirmation-stage numbers are reported as results.
\item \textbf{Matched tuning budget.} Every arm receives the same number of
      trials: $T = 10$ in the graph-only campaigns and $T = 40$ in the wide
      campaign of \S\ref{sec:wide}. The cosine baseline has a smaller native
      search space, so we expand its grid to keep the budget --- and hence the
      selection bias --- equal. Equal trials is not equal \emph{coverage} of
      the space, which we quantify and treat as a caveat against the
      higher-dimensional arms.
\item \textbf{Paired analysis.} All arms share the same seed set, so we
      difference per seed against the baseline. This cancels the dominant
      variance component and tightens the standard error on STS-B by
      roughly $5\times$ ($0.2 \rightarrow 0.036$).
\item \textbf{Distribution-free testing.} We report an exact two-sided sign
      test alongside the paired $t$-test, because accuracy on small
      evaluation sets is quantised (e.g.\ MRPC accuracy moves in steps of
      $0.245$ on 408 examples) and the parametric standard error is
      misleadingly small.
\item \textbf{Graph-health assertions.} Every run logs realised edge density.
      Runs whose graph is empty or complete are flagged, not silently averaged.
\item \textbf{Controls.} Every comparison includes a \texttt{no\_graph} arm
      (self-loops only), and every new normalisation knob is additionally run
      on a backbone that should be unaffected.
\item \textbf{Warm caches.} Hidden states are precomputed before any arm runs.
      \TODO{cross-ref to \S\ref{sec:artifacts} cache/RNG confound}
\end{enumerate}

\paragraph{Setup.} Frozen backbone, 2 epochs, Adam with learning rate
$2\times10^{-4}$, weight decay $0$, batch size $32$; two \texttt{GATConv}
layers with hidden width $128$, jumping-knowledge concatenation, projection
dimension $256$. \TODO{confirm these match the released \textsc{Glot} config
line by line and state any deviation}

\paragraph{Tasks and metrics.} CoLA (MCC), STS-B (Spearman), MRPC (F1),
RTE (accuracy). \TODO{decide whether to add MTEB / remaining GLUE tasks}

\paragraph{Backbones.} BERT-base, RoBERTa-base, ModernBERT-base,
SmolLM2-360M and TinyLlama-1.1B-Chat.
\TODO{state which backbones appear in which table once all runs land}

%-------------------------------------------------------------------------
\section{Reproduction of \textsc{Glot}}
\label{sec:repro}

Before making any claim about \textsc{Glot} we establish how closely we
reproduce it. Table~\ref{tab:repro} compares our re-run against the published
BERT row.

\begin{table}[t]
\caption{Reproduction against \citet{mantri2026glot}, Table 1, BERT backbone.
Our numbers are confirmation-stage means over 15 seeds; the published numbers
are single-seed maxima over a hyperparameter grid. \S\ref{sec:estimator} shows
that this difference of \emph{estimator}, not of implementation, accounts for
the gap on three of the four tasks.}
\label{tab:repro}
\begin{center}
\begin{tabular}{lcccc}
\toprule
 & CoLA & STS-B & MRPC & RTE \\
 & (MCC) & (Spearman) & (F1) & (Acc) \\
\midrule
Published \textsc{Glot}      & 47.49 & 83.86 & 82.58 & 59.21 \\
Ours (baseline, mean of 15)  & 47.01 & 82.61 & 81.93 & 55.72 \\
\midrule
$\Delta$                     & $-0.48$ & $-1.25$ & $-0.65$ & $-3.49$ \\
\bottomrule
\end{tabular}
\end{center}
\end{table}

\subsection{The published numbers are single-draw grid maxima}
\label{sec:estimator}

\citet{mantri2026glot} state that ``for all experiments, we used a fixed random
seed of 42''. Every published quality number is therefore a single draw. Their
Table~8 additionally reports the best $\tau$ per task, selected on the GLUE
development split --- which is also the split they report. The only dispersion
reported anywhere in that paper is on batch \emph{runtime} (their Table~5,
``mean $\pm$ standard deviation over 10 measurements''); no quality metric
carries an interval.

We report means over 15 seeds. These are not the same quantity.
Table~\ref{tab:estimator} re-scores \emph{our own} baseline runs under both
estimators.

\begin{table}[t]
\caption{The same baseline runs, summarised three ways. The published number is
a maximum over a search at one seed; the row we report is a mean over seeds.
Read down the CoLA column: the published $47.49$ lies between our per-seed
maximum and our per-configuration maximum, and $1.63$ seed-standard-deviations
above our mean.}
\label{tab:estimator}
\begin{center}
\begin{tabular}{lcccc}
\toprule
 & \multicolumn{2}{c}{CoLA (MCC)} & \multicolumn{2}{c}{STS-B (Spearman)} \\
\cmidrule(lr){2-3}\cmidrule(lr){4-5}
Estimator & value & vs.\ published & value & vs.\ published \\
\midrule
Published \textsc{Glot}          & 47.49 & --- & 83.86 & --- \\
\midrule
Mean over 15 seeds \emph{(ours)} & 45.37 & $-2.12$ & 82.53 & $-1.33$ \\
Max over 15 seeds                & 48.17 & $\mathbf{+0.68}$ & 83.11 & $-0.75$ \\
Max over 40 tuning configs       & 50.16 & $\mathbf{+2.67}$ & 82.95 & $-0.91$ \\
\midrule
Seed std / spread                & \multicolumn{2}{c}{$1.30$ / $5.46$} &
                                   \multicolumn{2}{c}{$0.33$ / $1.15$} \\
\bottomrule
\end{tabular}
\end{center}
\end{table}

\subsection{Matched comparison: their seed, their grid}

The comparison that removes the estimator difference fixes seed $42$ and sweeps
the published Table~6 grid, $78$ configurations per task
(Table~\ref{tab:seed42}).

\begin{table}[t]
\caption{Best score at \textsc{Glot}'s own seed ($42$) over \textsc{Glot}'s own
Table~6 grid, $78$ configurations per task. CoLA and RTE reproduce and are
exceeded; only STS-B retains a residual.}
\label{tab:seed42}
\begin{center}
\begin{tabular}{lccl}
\toprule
Task & ours @ seed 42 & published & selected configuration \\
\midrule
CoLA  & \textbf{50.52} & 47.49 & $\tau{=}0.3$, $K{=}4$, $h{=}256$, lr $10^{-3}$, jk\,=\,max \\
RTE   & \textbf{59.57} & 59.21 & $\tau{=}0.1$, $K{=}4$, $h{=}128$, lr $10^{-3}$, jk\,=\,max \\
STS-B & 82.99          & 83.86 & $\tau{=}0.3$, $K{=}4$, $h{=}64$, lr $2\times10^{-4}$, jk\,=\,cat \\
\bottomrule
\end{tabular}
\end{center}
\end{table}

Two of the three tasks reproduce. The STS-B residual is stable at roughly
$0.8$ under every estimator we tried and we do not reach $83.86$ in $495$ runs;
we leave it unexplained rather than attribute it. The most likely cause is a
configuration ambiguity we cannot resolve from the artifacts: the released
README specifies 3 epochs, hidden width $256$ and $\tau=0.8$, whereas the paper
specifies 2 epochs, hidden width $128$ and $\tau=0.6$. We followed the paper.

\subsection{The published \texorpdfstring{$\tau$}{tau} curve is not separable
from seed noise}
\label{sec:tau-noise}

Table~8 of \citet{mantri2026glot} is read as a $\tau$-sensitivity result. At
the noise level we measure, it cannot support that reading
(Table~\ref{tab:tau-noise}).

\begin{table}[t]
\caption{Reported $\tau$ sensitivity against measured variation at
\emph{fixed} $\tau$. Both published rows are flat with a single spike --- $+7.7$
above its neighbours on CoLA, $+6.5$ on RTE --- and in each case that spike is
the value promoted to their Table~1.}
\label{tab:tau-noise}
\begin{center}
\begin{tabular}{lc}
\toprule
Source of variation & range \\
\midrule
Published CoLA $\tau$ sweep ($39.62 \rightarrow 47.49$) & 7.87 \\
Published RTE $\tau$ sweep ($50.90 \rightarrow 59.21$)  & 8.31 \\
\midrule
Our CoLA spread across 15 seeds at \emph{fixed} $\tau$   & 5.46 \\
Our RTE spread across configurations at \emph{fixed} seed 42 & 11.92 \\
\bottomrule
\end{tabular}
\end{center}
\end{table}

\paragraph{Consequence for this paper.} Because the published numbers are grid
maxima and ours are seed means, we never compare against them directly. Every
comparison in the remainder of the paper is against \emph{our} baseline, run
under identical conditions, with the same seeds and the same tuning budget.

%-------------------------------------------------------------------------
\section{The graph is not calibrated across backbones}
\label{sec:calibration}

\subsection{Edge density under the published $\tau$ grid}

Table~\ref{tab:density} reports the realised edge density of
Eq.~\ref{eq:glot-adj} at each $\tau$ the original work searched, together with
the underlying token-cosine distribution.

\begin{table}[t]
\caption{Edge density induced by \textsc{Glot}'s own search space
$\tau \in \{0.1, 0.3, 0.6\}$, with the token-cosine distribution that produces
it. A density near $1$ means the GNN receives a near-complete graph and
therefore no selective relational structure.}
\label{tab:density}
\begin{center}
\begin{tabular}{lcccccc}
\toprule
\multirow{2}{*}{Backbone} & \multicolumn{3}{c}{$\cos$ percentile} &
\multicolumn{3}{c}{density at $\tau =$} \\
\cmidrule(lr){2-4}\cmidrule(lr){5-7}
 & p10 & p50 & p90 & $0.1$ & $0.3$ & $0.6$ \\
\midrule
BERT-base        & 0.098 & 0.328 & 0.595 & 0.876 & 0.505 & \textbf{0.103} \\
RoBERTa-base     & 0.693 & 0.821 & 0.897 & 1.000 & 1.000 & \textbf{0.991} \\
ModernBERT-base  & 0.076 & 0.648 & 0.768 & 0.866 & 0.759 & \textbf{0.599} \\
SmolLM2-360M     & 0.322 & 0.774 & 0.911 & 0.956 & 0.908 & 0.741 \\
TinyLlama-1.1B   & 0.016 & 0.343 & 0.588 & 0.762 & 0.533 & \textbf{0.073} \\
\bottomrule
\end{tabular}
\end{center}
\end{table}

Two backbones sit in a usable regime (BERT, TinyLlama); three do not. The
mechanism is visible in the percentile columns: RoBERTa's \emph{tenth}
percentile token cosine is $0.701$, which exceeds the largest $\tau$ in the
grid. There is therefore no setting within the published search space that
makes RoBERTa's token graph sparse. Whatever the RoBERTa results in
\citet{mantri2026glot} measure, it is not the benefit of a selectively
constructed graph.

\subsection{Does the degenerate graph change conclusions?}

\TODO{This is the load-bearing experiment for the section and is still running.
Fill Table~\ref{tab:roberta} with RoBERTa under (i) the published absolute
$\tau$, (ii) density-matched \texttt{tau\_quantile}, (iii) Stage A, and
(iv) \texttt{no\_graph}. The question to answer explicitly: once RoBERTa's
graph is made sparse, does performance improve, and does \texttt{no\_graph}
still match the baseline?}

\begin{table}[t]
\caption{RoBERTa-base with a degenerate versus density-matched graph.}
\label{tab:roberta}
\begin{center}
\begin{tabular}{lcccc}
\toprule
Arm & CoLA & STS-B & MRPC & RTE \\
\midrule
published $\tau$ (density $0.992$) & \PEND & \PEND & \PEND & \PEND \\
density-matched                    & \PEND & \PEND & \PEND & \PEND \\
Stage A                            & \PEND & \PEND & \PEND & \PEND \\
\texttt{no\_graph}                 & \PEND & \PEND & \PEND & \PEND \\
\bottomrule
\end{tabular}
\end{center}
\end{table}

%-------------------------------------------------------------------------
\section{A two-part correction}
\label{sec:fix}

\subsection{Relative-density thresholding}
We replace the absolute threshold with a quantile criterion: retain the top
$q$ fraction of candidate pairs by similarity, so that realised density is
$q$ by construction and comparable across backbones.

\subsection{Median-based input scaling}
\label{sec:scaling}
Density alone is not sufficient. Some backbones carry extreme token-norm
outliers --- ``attention sinks'' --- which distort any mean-based
normalisation. For ModernBERT the ratio of mean to median token norm is
$7.66$ at layer 12 and $9.87$ at layer 22, with $\max/\text{median} \approx
100$. Normalising by the mean therefore places the \emph{typical} token far
below the target scale. We normalise by the median instead.

\subsection{Result}
Table~\ref{tab:modernbert} varies the two corrections independently.

\begin{table}[t]
\caption{ModernBERT-base on CoLA (MCC). The two corrections are necessary
jointly: density-matching alone is actively harmful. The BERT control confirms
the knob is inert where it should be --- forcing the correction on a backbone
whose graph is already well-conditioned costs $2.13$ MCC, as expected.}
\label{tab:modernbert}
\begin{center}
\begin{tabular}{llcc}
\toprule
Backbone & Configuration & Layer & MCC \\
\midrule
\multirow{5}{*}{ModernBERT-base}
 & $\tau = 0.6$ (published)                    & 12 & 21.71 \\
 & \texttt{tau\_quantile} $= 0.05$             & 12 & \phantom{0}9.16 \\
 & \texttt{tau\_quantile} $+$ \texttt{rms}     & 12 & 23.15 \\
 & \texttt{tau\_quantile} $+$ \texttt{median}  & 12 & \textbf{27.15} \\
 & $\tau = 0.6$ $+$ \texttt{median}            & 12 & \TODO{missing cell} \\
\midrule
\multirow{3}{*}{ModernBERT-base}
 & $\tau = 0.6$ (published)                    & \phantom{0}4 & 13.18 \\
 & \texttt{tau\_quantile} $= 0.05$             & \phantom{0}4 & 16.91 \\
 & \texttt{tau\_quantile} $= 0.10$             & \phantom{0}4 & 12.24 \\
\midrule
\multirow{2}{*}{BERT-base \emph{(control)}}
 & $\tau = 0.6$ (published)                    & 12 & 45.54 \\
 & \texttt{tau\_quantile} $= 0.05$             & 12 & 43.41 \\
\bottomrule
\end{tabular}
\end{center}
\end{table}

\paragraph{An honest gap.} The cell $\tau = 0.6 + \texttt{median}$ was not run,
so the $+5.44$ improvement cannot yet be decomposed into a density
contribution and a scale contribution. We can currently claim only that the
two corrections \emph{together} help and that density-matching alone does not.
\TODO{run the missing cell and decompose}

%-------------------------------------------------------------------------
\section{Does hyperbolic geometry help?}
\label{sec:hyperbolic}

\subsection{BERT}
Table~\ref{tab:bert-arms} reports all arms on BERT under the \emph{narrow}
graph-only search. It is superseded by \S\ref{sec:wide} and is retained only
because the contrast between the two is itself the result.

\begin{table}[t]
\caption{All arms on BERT-base under the \textbf{narrow, graph-only search},
confirmation stage. \textbf{This table is superseded by
Table~\ref{tab:wide}}: only the graph-construction knobs were searched here,
while the learning rate, depth, width, jumping-knowledge mode and projection
width were inherited from values the original authors chose for a
\emph{Euclidean} pooler. That asymmetry is what produced the Stage A result we
retract in \S\ref{sec:stsb-retracted}, and it affects every row. Sample sizes
differ by row and are marked; at $n=3$ the minimum detectable effect is $0.64$
(STS-B), $0.68$ (MRPC) and $6.12$ (RTE), so no row in the lower block can
support a directional claim.}
\label{tab:bert-arms}
\begin{center}
\begin{tabular}{lccc}
\toprule
Arm & STS-B & MRPC & RTE \\
\midrule
baseline (cosine) \hfill($n{=}15$) & 82.61 & 81.93 & 55.72 \\
A \hfill($n{=}15$)                 & \textbf{82.83} & 81.97 & \textbf{55.84} \\
\midrule
\texttt{no\_graph} \hfill($n{=}3$) & 82.48 & 81.82 & 55.47 \\
B    \hfill($n{=}3$) & 82.22 & 81.41 & 54.75 \\
C    \hfill($n{=}3$) & 82.79 & 81.61 & 54.27 \\
AB   \hfill($n{=}3$) & 81.93 & 81.37 & 55.72 \\
AC   \hfill($n{=}3$) & 82.81 & 81.62 & 54.39 \\
BC   \hfill($n{=}3$) & 82.03 & 81.22 & 54.39 \\
ABC  \hfill($n{=}3$) & 82.04 & \textbf{82.21} & 54.63 \\
\bottomrule
\end{tabular}
\end{center}
\end{table}

\paragraph{Reconciling with Table~\ref{tab:wide}.} In this table
\texttt{no\_graph} sits \emph{below} the baseline on all three tasks
($-0.13$, $-0.11$, $-0.25$); in the wide search it sits \emph{above} it on both
tasks measured there ($+0.035$ STS-B, $+0.833$ CoLA). The apparent
contradiction is not one. On the task common to both, STS-B, the two estimates
are $-0.13$ (here, $n=3$, unpaired) and $+0.035$ (Table~\ref{tab:wide},
$n=15$, paired, $7$ seeds positive and $8$ negative, sign-test $p = 1.00$).
Both are far inside the noise: the minimum detectable effect is $0.64$ at
$n=3$ and $0.225$ at $n=15$. The sign of a difference this small is not
information, and we draw no conclusion from either number in isolation.

What we do claim is narrower and rests on Table~\ref{tab:wide} alone, where
every arm received an equal budget over the full grid: \texttt{no\_graph} is
statistically indistinguishable from the baseline on STS-B and ahead of it on
CoLA, and no arm is distinguishable from \texttt{no\_graph} in the direction
that would show the graph carrying information (\S\ref{sec:vs-nograph}).
Readers should treat Table~\ref{tab:bert-arms} as a record of the superseded
protocol, not as evidence.

\subsection{The one positive result}
\label{sec:stsb-retracted}
\paragraph{Retracted.} Under the graph-only search this was the paper's
headline: Stage A on STS-B gave $\bar{d} = +0.223$, $t(14) = 6.20$, sign-test
$p = 6\times10^{-5}$, $15/15$ seeds positive, and replicated on 12 held-out
seeds. It does not survive \S\ref{sec:wide}. Once the Euclidean baseline is
allowed to search the same depth/width grid as the hyperbolic arms, the same
comparison gives $\bar{d} = +0.028$, $t = 0.69$, $9/5$ seeds --- a clean null.

The effect was real but its \emph{cause} was not geometry: the baseline had been
pinned at $K=2$, $h=128$ while nothing else was, and it recovers the entire gap
once it can choose $K=4$. We report this prominently rather than quietly
dropping it, because it is the clearest illustration of this paper's own thesis
--- that a pooling comparison is only as trustworthy as the budget given to its
control.

\subsection{A false positive we retracted}
An earlier $n=3$ analysis showed Stage A ahead on MRPC and we initially
recorded it as a second positive. At $n=15$ the effect is a clean null:
$\bar{d} = +0.035$, $t = 0.40$, $7$ seeds positive and $7$ negative,
sign-test $p = 1.00$. We report this because the $n=3$ minimum detectable
effect ($0.68$) was larger than any effect we could plausibly have measured;
the apparent result was a property of the sample size, not of the method.

\subsection{Widening the search to the paper's full grid}
\label{sec:wide}

Every result above searched only graph-construction knobs, inheriting the
learning rate, depth, width, jumping-knowledge mode and projection width from
the values the original authors selected --- for a \emph{Euclidean} pooler. That
is a confound, and it points against the hyperbolic arms. We therefore repeat
the campaign with the paper's own Table~6 optimizer/architecture grid added on
top of each arm's graph grid, applied identically to every arm including the
baseline: 40 trials per arm, 15 confirmation seeds, 990 runs.

\paragraph{The learning-rate hypothesis is false.} We expected Stages B and C to
have been trained at the wrong learning rate. They were not. \emph{Every arm,
including the baseline, selected $\mathrm{lr}=2\times10^{-4}$} --- exactly the
pinned value. What moved instead was depth and width: most arms preferred $K=4$
over the pinned $K=2$, and hidden width ranged over $\{64,128,256\}$.

\begin{table}[t]
\caption{BERT-base, wide search, confirmation stage, $n=15$ paired seeds.
Verdicts require the paired $t$-test \emph{and} the exact sign test to agree;
Bonferroni is over the 8 arms tested against a common baseline.}
\label{tab:wide}
\begin{center}
\begin{tabular}{lccccl}
\toprule
Arm & $\Delta$ STS-B & sign $p$ & $\Delta$ CoLA & sign $p$ & verdict \\
\midrule
C   & $+0.195$ & 0.035 & $\mathbf{+1.416}$ & 0.0074 & positive, fails Bonferroni \\
AC  & $+0.239$ & 0.118 & $+0.448$ & 0.302 & n.s. \\
\texttt{no\_graph} & $+0.035$ & 1.000 & $+0.833$ & 0.118 & n.s. \\
A   & $+0.028$ & 0.424 & $+0.700$ & 0.118 & n.s. \\
\midrule
AB  & $-0.747$ & 0.00006 & $-1.236$ & 0.118 & \textbf{harmful} (both) \\
ABC & $-0.837$ & 0.00098 & $-0.528$ & 0.607 & \textbf{harmful} (STS-B) \\
B   & $-1.065$ & 0.00006 & $-1.263$ & 0.302 & \textbf{harmful} (both) \\
BC  & $-1.242$ & 0.00006 & $+0.224$ & 1.000 & \textbf{harmful} (STS-B), n.s.\ CoLA \\
\bottomrule
\end{tabular}
\end{center}
\end{table}

\paragraph{Every arm that beats the baseline somewhere.} Five of the eight arms
have a positive point estimate against the baseline on at least one task:
C ($+0.195$ STS-B, $+1.416$ CoLA), AC ($+0.239$, $+0.448$), \texttt{no\_graph}
($+0.035$, $+0.833$), A ($+0.028$, $+0.700$), and BC ($+0.224$ on CoLA only,
against $-1.242$ on STS-B). We state this explicitly because a table grouped by
verdict can hide a per-task win --- BC's CoLA gain sits inside the block we
label harmful.

Two cautions apply to that list, and they are the reason we do not promote it
to a headline. First, ``beats the baseline on at least one of two tasks'' is
itself a selection over tasks, and is the same maximisation this paper
criticises in \S\ref{sec:estimator}; of the five arms, only C, AC, A and
\texttt{no\_graph} are positive on \emph{both}, and none of the four is
significant on both tests after correction. Second, four of those five --- every
one except BC --- are also matched or beaten by \texttt{no\_graph}
(\S\ref{sec:vs-nograph}), so their advantage over \textsc{Glot}'s baseline is
not evidence for the geometry they add. We therefore report the full arm
$\times$ task grid in every table rather than a filtered list of winners, and
ask the reader to apply the \texttt{no\_graph} column before reading any row as
a success.

\paragraph{MRPC and RTE are unresolved.} The wide campaign covers CoLA and
STS-B only. On MRPC and RTE we have the full arm set at $n=3$
(Table~\ref{tab:bert-arms}), where ABC leads on MRPC ($82.21$ vs.\ $81.93$) and
A leads on RTE ($55.84$ vs.\ $55.72$). Both gaps are far below the minimum
detectable effect at that sample size ($0.68$ and $6.12$ respectively), so
neither can be called a win or a loss. Until those two tasks are run at $n=15$
under the wide grid, any statement of the form ``no arm beats the baseline on
any task'' is stronger than our evidence supports, and we do not make it.
\TODO{run MRPC and RTE at $n=15$ under the wide grid and fold into
Table~\ref{tab:wide}}

\paragraph{Stage B is harmful; Stage C is not.} The earlier framing --- ``the
hyperbolic readout and GNN are both harmful'' --- conflated two components that
move in opposite directions. Every arm containing the hyperbolic \emph{readout}
(B) is negative on STS-B, all four survive Bonferroni, and three of them have
$0/15$ or $1/14$ seeds positive. Every arm containing the hyperbolic \emph{GNN}
without the readout (C, AC) is positive on both tasks. Stage C alone reaches
$+1.416$ MCC on CoLA with $13/15$ seeds ($p=0.0074$), which is significant on
both tests but does not survive Bonferroni across 8 arms. We therefore report
Stage C as \emph{suggestive and consistent in direction across two tasks},
not as established.

\paragraph{Stage A is a null.} With the baseline free to choose its own depth
and width, Stage A's STS-B advantage disappears entirely
(\S\ref{sec:stsb-retracted}).

\paragraph{\texttt{no\_graph} survives everything.} Removing the token graph is
statistically indistinguishable from the full cosine graph on both tasks here
($+0.035$ and $+0.833$), as it was on the decoder (\S\ref{sec:decoder}). Across
every backbone, task and search budget we have tried, we have not found a
setting in which the constructed graph demonstrably carries information.

\TODO{Coverage caveat for the appendix: 40 trials cover 6.2\% of
\texttt{no\_graph}'s space but $<0.01\%$ of ABC's. Equal \emph{trials} is not
equal \emph{coverage}, and this asymmetry disfavours the complex arms.}

\subsection{Differencing against the no-graph control}
\label{sec:vs-nograph}

Every comparison above is referenced to the cosine baseline. But
\texttt{no\_graph} \emph{also} beats that baseline on CoLA ($+0.833$), so an
arm's advantage over the baseline may be no more than ``not using
\textsc{Glot}'s graph''. The quantity that isolates geometry is the difference
against \texttt{no\_graph}, reported in Table~\ref{tab:vs-nograph}.

\begin{table}[t]
\caption{BERT-base, wide search, $n=15$ paired seeds, referenced to the
\texttt{no\_graph} control rather than to the cosine baseline. Arms are ordered
by CoLA. The cosine baseline is \emph{below} the control on both tasks.}
\label{tab:vs-nograph}
\begin{center}
\begin{tabular}{lcccccc}
\toprule
 & \multicolumn{3}{c}{CoLA (MCC)} & \multicolumn{3}{c}{STS-B (Spearman)} \\
\cmidrule(lr){2-4}\cmidrule(lr){5-7}
Arm & mean & $\Delta$ & pos/neg & mean & $\Delta$ & pos/neg \\
\midrule
C                  & \textbf{46.78} & $+0.583$ & 8/7   & 82.73 & $+0.160$ & 12/3 \\
\texttt{no\_graph} & 46.20 & \emph{ref} & ---         & 82.57 & \emph{ref} & --- \\
A                  & 46.07 & $-0.133$ & 4/11          & 82.56 & $-0.007$ & \phantom{0}7/8 \\
AC                 & 45.82 & $-0.385$ & 6/9           & \textbf{82.77} & $+0.203$ & 10/5 \\
BC                 & 45.59 & $-0.609$ & 4/11          & 81.29 & $-1.277$ & \phantom{0}0/15 \\
baseline           & 45.37 & $-0.833$ & 4/11          & 82.53 & $-0.035$ & \phantom{0}8/7 \\
ABC                & 44.84 & $-1.361$ & 5/10          & 81.70 & $-0.872$ & \phantom{0}1/14 \\
AB                 & 44.13 & $-2.069$ & 3/12          & 81.79 & $-0.782$ & \phantom{0}1/14 \\
B                  & 44.11 & $-2.095$ & 3/12          & 81.47 & $-1.100$ & \phantom{0}0/15 \\
\bottomrule
\end{tabular}
\end{center}
\end{table}

\paragraph{Stage C's advantage does not survive the change of reference.} Of
its $+1.416$ MCC over the baseline, $+0.833$ is shared with the control that
deletes the graph entirely. The residual $+0.583$ has $8$ seeds positive and
$7$ negative --- an exact sign-test $p$ of $1.00$, which is as null as a result
can be. On STS-B the residual is $+0.160$ with $12/3$ seeds, significant on the
sign test but not on the $t$-test.

\paragraph{What we are willing to claim.} We do not claim that hyperbolic
message passing improves this architecture. The defensible statements are
narrower and all three are supported on both tasks: the hyperbolic readout
(Stage B) is reliably harmful; the hyperbolic GNN (Stage C) is never harmful
and is directionally positive; and neither is distinguishable from simply
removing the token graph. The most economical description of
Table~\ref{tab:vs-nograph} is that \textsc{Glot}'s cosine graph is slightly
worse than no graph, and that the geometric variants differ from each other
mainly in how much damage the readout does.

\subsection{Decoder backbone}
\label{sec:decoder}
We repeat the full arm sweep on TinyLlama-1.1B, selected over SmolLM2-360M on
measured geometry rather than preference: TinyLlama sits in BERT's density
regime ($0.073$ at $\tau=0.6$, against BERT's $0.103$) whereas SmolLM2 does not
($0.741$; Table~\ref{tab:density}).

\begin{table}[t]
\caption{TinyLlama-1.1B on STS-B (Spearman), confirmation stage, $n=15$ shared
seeds. $\Delta$ is the paired per-seed difference against the baseline;
``pos/neg'' counts seeds in each direction and $p$ is an exact two-sided sign
test. The minimum detectable effect at this sample size is $0.503$.
Stage A does \emph{not} transfer from BERT; removing the graph entirely costs
nothing; and Stages B and C are significantly harmful.}
\label{tab:decoder}
\begin{center}
\begin{tabular}{lccccccc}
\toprule
Arm & mean & std & density & $\Delta$ & $t$ & sign $p$ & pos/neg \\
\midrule
AC                 & 80.15 & 0.54 & 0.122 & $+0.205$ & $1.01$ & $0.302$ & 10/5 \\
A                  & 79.99 & 0.51 & 0.194 & $+0.043$ & $0.34$ & $0.607$ & \phantom{0}9/6 \\
\texttt{no\_graph} & 79.95 & 0.48 & 0.095 & $+0.001$ & $0.01$ & $1.000$ & \phantom{0}8/7 \\
baseline           & 79.95 & 0.49 & 0.121 & \emph{ref} & --- & --- & --- \\
\midrule
BC                 & 78.78 & 0.95 & 0.196 & $-1.165$ & $-4.97$ & $0.00098$ & \phantom{0}1/14 \\
AB                 & 76.39 & 1.87 & 0.195 & $-3.556$ & $-6.95$ & $0.00006$ & \phantom{0}0/15 \\
ABC                & 75.44 & 3.58 & 0.195 & $-4.506$ & $-4.85$ & $0.00006$ & \phantom{0}0/15 \\
\bottomrule
\end{tabular}
\end{center}
\end{table}

\paragraph{Stage A does not transfer.} The STS-B effect that is highly
significant on BERT ($+0.223$, $15/15$ seeds) is absent here: $+0.043$ with
$9/6$ seeds and $p = 0.61$. This is not merely a power failure --- Stage A's
paired standard error on this backbone is $0.127$, so an effect the size of the
BERT one would have been detectable. We therefore report the STS-B result as
\emph{backbone-specific} rather than as a property of hyperbolic token graphs.

\paragraph{Removing the graph costs nothing.} The \texttt{no\_graph} arm is
statistically indistinguishable from the full cosine graph: $\Delta = +0.001$,
$t = 0.01$, sign-test $p = 1.00$, $8$ seeds positive and $7$ negative. This is
not a calibration artifact --- TinyLlama's baseline density is $0.121$, squarely
in the regime the method is designed for (Table~\ref{tab:density}). On this
backbone and task the constructed token graph carries no measurable
information.

\paragraph{Stages B and C are actively harmful.} Unlike every other negative in
this paper, these are significant: ABC loses $4.51$ points with $0/15$ seeds
positive, AB loses $3.56$ with $0/15$, and BC loses $1.17$ with $1/14$. The
variance ordering is monotone in the number of hyperbolic components
($0.48 \rightarrow 0.95 \rightarrow 1.87 \rightarrow 3.58$), consistent with the
readout and message-passing operators acting as variance amplifiers rather than
as useful inductive bias.

\paragraph{Confound to resolve.} Realised density spans $0.095$--$0.196$ across
arms ($2.1\times$), so differences could in principle reflect sparsity rather
than geometry. This matters least for the two arms that matter most: AC, the
largest positive, sits at density $0.122$ against the baseline's $0.121$, and
\texttt{no\_graph} is a null regardless. \TODO{add a density-matched Euclidean
control at each arm's realised density to close this off}

In the tuning stage Stage A led the baseline by $+0.16$
($81.51$ vs $81.35$) --- against a selection-bias inflation of $0.83$ at this
arm's noise level (\S\ref{sec:artifacts}), i.e.\ within what the search alone
produces. Its confirmation mean is $1.52$ below its tuning maximum.

\subsection{Structural (non-geometric) graphs}
Since CoLA, RTE and MRPC are flat for every geometric arm, we also test whether
a \emph{structural} prior helps where a similarity prior does not: positional
edges connecting tokens with $|i - j| \leq w$, positional-only graphs, their
union with Stage A, and $k$-NN graphs.

\begin{table}[t]
\caption{Structural arms on CoLA, tuning stage. These are maxima over 10
trials and are \emph{not} results: at CoLA's confirmation-stage noise level
($\sigma = 2.09$) the expected inflation is $\approx 3.22$, which exceeds every
gap in the table.}
\label{tab:structural}
\begin{center}
\begin{tabular}{lccc}
\toprule
Arm & trials & best & mean \\
\midrule
baseline           & 10 & 45.32 & 43.42 \\
POS                & 10 & 47.08 & 44.50 \\
POS\_ONLY          & \phantom{0}4 & 46.80 & 44.38 \\
A\_POS             & 10 & 44.87 & 43.48 \\
KNN                & \phantom{0}4 & 45.32 & 44.84 \\
\bottomrule
\end{tabular}
\end{center}
\end{table}

\TODO{confirmation stage for structural arms on CoLA, RTE, MRPC}

%-------------------------------------------------------------------------
\section{Two measurement artifacts}
\label{sec:artifacts}

\subsection{$\delta$-hyperbolicity on raw activations measures attention sinks}
A natural way to motivate hyperbolic pooling is to show that token
representations are already tree-like, via Gromov's $\delta$-hyperbolicity.
Doing so on raw activations is misleading. ModernBERT at layer 16 has
$\delta_{\text{eucl}} = 0.0021$, which would indicate an almost perfectly
tree-like space. Its ratio of maximum to median token norm is $156$. Computing
$\delta$ on the \emph{angular} metric $\arccos(\cos(\cdot,\cdot))$, which is
invariant to these outliers, gives $0.2127$ --- statistically indistinguishable
from BERT's $0.2114$. Flan-T5 shows the same pattern more extremely
($357\times$ norm outliers).

We therefore recommend reporting angular $\delta$ alongside Euclidean $\delta$,
together with the norm-outlier ratio, in any work that motivates hyperbolic
geometry from measured hyperbolicity. \TODO{table of $\delta_{\text{eucl}}$,
$\delta_{\text{ang}}$ and $\max/\text{median}$ norm for all backbones and
layers}

\subsection{Selection bias makes tuning maxima unusable}
With $T$ trials and per-seed noise $\sigma$, the expected maximum over trials
is inflated even for an arm identical to the baseline. In our setting the
inflation is $\approx 1.83$ (STS-B on BERT, $\sigma = 1.19$) and $\approx 0.83$
(STS-B on TinyLlama, $\sigma = 0.54$). Several apparent leads in this project
were smaller than this and did not survive confirmation.

A related trap: we initially assumed CoLA's confirmation noise was
$\sigma = 0.81$ based on a layer-12 probe. The true value at layer 8 is
$2.09$. Any tuning-stage lead must be compared against the noise floor
\emph{of the configuration being confirmed}, measured directly.

\subsection{Other pitfalls}
\TODO{write up, briefly: (i) cold vs.\ warm feature cache consumes RNG draws
via the shuffled loader and shifts classifier initialisation, worth
$\approx 5$ MCC on CoLA and biasing whichever arm runs first; (ii) selecting a
backbone layer on a single task overfits --- layer 8 beat layer 12 by
$+4.62$ MCC on CoLA but failed to transfer ($+0.25$ STS-B, $-0.36$ RTE,
$-0.46$ MRPC); (iii) \texttt{edge\_dim} is nearly inert in this architecture.}

%-------------------------------------------------------------------------
\section{Related work}
\label{sec:related}

\TODO{Write this section. Required coverage, with citations to be fetched and
verified programmatically --- do not draft from memory:
(i) sentence representation and pooling;
(ii) graph-based and structure-aware pooling, including \textsc{Glot};
(iii) anisotropy and the geometry of contextual embeddings, which is the
direct cause of the calibration failure in \S\ref{sec:calibration};
(iv) attention sinks and massive activations, which cause the artifact in
\S\ref{sec:artifacts};
(v) hyperbolic representation learning and hyperbolic GNNs;
(vi) reproducibility and evaluation-protocol work on selection bias and seed
variance.}

%-------------------------------------------------------------------------
\section{Limitations}
\label{sec:limitations}

\begin{itemize}
\item One reproduction residual is unexplained. At the authors' seed and grid
      we exceed the published CoLA and RTE numbers (Table~\ref{tab:seed42}),
      but STS-B remains $0.87$ short and stays short under every estimator we
      tried. We suspect the README/paper configuration discrepancy (3 vs.\ 2
      epochs, hidden $256$ vs.\ $128$) but have not confirmed it.
\item Our own reported numbers select the epoch on the same GLUE development
      split we report, because GLUE test labels are not public. This is the
      practice we criticise in \S\ref{sec:estimator}, applied at smaller
      scale: we select over epochs and configurations, they additionally
      select over $\tau$ at one seed. A split of the development set into
      selection and reporting halves would remove the objection entirely and
      we regard it as the most valuable single improvement to this protocol.
\item Several arms are reported at $n = 3$ and are underpowered by our own
      criterion; they are included for ordering only.
\item The correction in \S\ref{sec:fix} is validated on one backbone and one
      task, and its two components are not yet independently identified.
\item Our negative result concerns the specific hyperbolic instantiations in
      \S\ref{sec:stages}. It does not rule out other geometric formulations,
      and in particular structured layers are known to require structure-aware
      learning rates and initialisation before their benefits appear
      \TODO{cite; this is the strongest objection to our negative result and
      must be addressed, not merely acknowledged}.
\item \TODO{state compute budget and total GPU-hours}
\end{itemize}

%-------------------------------------------------------------------------
\section{Conclusion}
\label{sec:conclusion}

We set out to test whether hyperbolic geometry improves graph-based sentence
pooling, and found that the prior question --- whether the graph helps at all
--- was not settled.

Three results stand. First, \textsc{Glot}'s graph construction has an unstated
calibration dependency: an absolute cosine threshold tuned on BERT produces
near-complete graphs on RoBERTa, ModernBERT and SmolLM2, and for RoBERTa no
value in the published search space produces a sparse graph at all. A two-part
correction --- relative-density thresholding plus median-based input scaling ---
recovers $+5.44$ MCC on ModernBERT/CoLA, and a BERT control confirms it is
inert where it should be.

Second, once every arm is given the same tuning budget over the paper's own
grid, the hyperbolic extensions separate cleanly rather than uniformly. The
hyperbolic readout is reliably harmful, losing up to $1.24$ Spearman with $0$ of
$15$ seeds positive and surviving Bonferroni correction. The hyperbolic GNN is
never harmful and is directionally positive on both tasks. But referenced
against a control that deletes the token graph, rather than against
\textsc{Glot}'s cosine baseline, even the hyperbolic GNN is a null. We do not
claim that hyperbolic geometry helps this architecture.

Third, and most consequentially for the family rather than for our extensions:
across five backbones, four tasks and roughly $2{,}000$ runs we never found a
setting in which the constructed token graph demonstrably carries information.
On CoLA, deleting it outright \emph{beats} the published construction. This is
a question about graph-based pooling in general, and it is not answered by
adding geometry to the pooler.

Underlying all three is a measurement point. The published results we set out
to improve on are single-seed maxima over a hyperparameter grid, evaluated on
the split they report; at the authors' own seed and grid we exceed their CoLA
number by $3.03$ MCC while our 15-seed mean sits $2.12$ below it. Both
statements are true of the same runs. A pooling comparison is only as
trustworthy as the budget given to its control and the estimator used to
summarise it --- a lesson this project learned twice at its own expense, in the
Stage A result we retracted in \S\ref{sec:stsb-retracted} and in the Stage C
result we decline to claim in \S\ref{sec:vs-nograph}.

\subsubsection*{Reproducibility statement}
\TODO{Point to the code release and the exact commit.}
All confirmation-stage results use the fixed seed set $\{1,\dots,15\}$, shared
across every arm so that comparisons are paired. Tuning uses a single held-out
seed disjoint from that set. Every run appends one row to a per-campaign CSV
recording the full configuration, the realised edge density, the per-epoch
metrics and the wall-clock time; these logs are released unmodified, including
runs that crashed or produced degenerate graphs. The analysis scripts that
produce every table in this paper read those CSVs directly and are released
alongside them. Hidden-state caches are built once, before any arm runs, with
\texttt{-{}-override\_precompute=1}; all reported runs then use
\texttt{-{}-override\_precompute=0}. This ordering is load-bearing rather than
an optimisation --- see \S\ref{sec:artifacts} --- and any attempt to reproduce
these numbers without it will shift results by several points on CoLA.

\subsubsection*{Acknowledgments}
\TODO{}

\bibliography{references}
\bibliographystyle{iclr2026_conference}

\appendix
\section{Appendix}

\subsection{Full hyperparameter search spaces}
\TODO{}

\subsection{Per-seed results}
\TODO{}

\subsection{Realised edge density for every reported run}
\TODO{}

\end{document}
